\section{Proof of pattern-consistent transfer}\label{supp:pattern89}
\subsection{Embedding and comparison with the earlier relaxation}
Fix an actual qualified parameter $\theta$ and its complete pattern $\sigma$. Give that pattern weight one and all others weight zero. Under the original law of its complete policy, set $m^\sigma_{tC}=\Pr(X_t\in C)$ and $z^\sigma_{tCD}=\Pr(X_t\in C,X_{t+1}\in D)$. Integrating the whole-cell transition bounds gives \eqref{eq:pattern-flow89}. The same occupation identity proved in Supplement~\ref{supp:occupation86} bounds the true cost difference between the two objectives in \eqref{eq:pattern-objective89}. This proves validity for the entire parameter class. There is no additional probability union over its members.

To compare the two relaxations, take any disaggregated feasible point and define
\begin{align*}
 y_{tCa}&=\sum_{\sigma:\sigma(t,C)=a}\lambda_\sigma,&
 q_{tCa}&=\sum_{\sigma:\sigma(t,C)=a}m^\sigma_{tC},\\
 z_{tCaD}&=\sum_{\sigma:\sigma(t,C)=a}z^\sigma_{tCD},&
 m_{tC}&=\sum_\sigma m^\sigma_{tC}.
\end{align*}
Conservation implies $\sum_Cm^\sigma_{tC}=\lambda_\sigma$ at every date, hence $0\le m^\sigma_{tC}\le\lambda_\sigma$ and $q_{tCa}\le y_{tCa}$. All assignment, flow and transition constraints of \eqref{eq:occupation-polytope86} follow by summation. Each retained pattern satisfies every certified conflict, so their convex combination satisfies its corresponding inequality in $y$. Both advantage objectives are unchanged by this projection. Consequently the disaggregated minimum is no smaller, and its maximum is no larger, than the earlier pair. Deleting only patterns certified impossible preserves coverage and shrinks the feasible set. Nonempty compactness follows from the actual-policy embedding, bounded total mass and closed linear constraints. This also proves existence of the displayed extrema and justifies ordinary linear-programming dual certificates.

\subsection{Quantitative comparison with an attained pattern}
Now assume every retained pattern is actual and has a qualification witness. Consider one pattern with positive weight and divide its flow variables by that weight. Its row sums define a time-inhomogeneous stochastic matrix $\widetilde P_t$ on the dated cells. At a row with zero mass choose any stochastic row within the same transition intervals; such a row exists because the original transition at any state in the nonempty cell supplies one.

Let $m_t$ be the true cell law of this pattern's original continuous-state policy, and $P_t$ its conditional cell-to-cell probabilities given the date-$t$ cell under that law. The law inside a cell can depend on the complete earlier history. Nevertheless, conditional expectation and the whole-cell hypothesis put every row of $P_t$ inside the same interval box. At an unvisited cell choose an arbitrary original-law row. This choice does not affect $m_tP_t=m_{t+1}$. No Markov assumption on the observed cell process has been imposed.

Two stochastic rows in this box differ in $\ell^1$ norm by at most $\delta_t$. Writing $\widetilde m_t$ for the normalized relaxed law, stochastic contraction gives
\begin{align*}
 \|\widetilde m_{t+1}-m_{t+1}\|_1
 &\le\|(\widetilde m_t-m_t)\widetilde P_t\|_1
       +\|m_t(\widetilde P_t-P_t)\|_1\\
 &\le\|\widetilde m_t-m_t\|_1+\delta_t.
\end{align*}
The initial laws agree. Since the distance between two probability vectors is at most two, induction bounds this distance by $e_t$.

Let $a_{tC}$ denote the true conditional expected reference advantage under this pattern at date $t$, with an arbitrary value in the enclosure on unvisited cells. Then $\ell_{tC\sigma(t,C)}\le a_{tC}\le u_{tC\sigma(t,C)}$. For the upper objective at that date,
\begin{align*}
 \sum_Cu_{tC\sigma(t,C)}\widetilde m_{tC}-\sum_Ca_{tC}m_{tC}
 &=\sum_Cu_{tC\sigma(t,C)}(\widetilde m_{tC}-m_{tC})\\
 &\quad+\sum_C(u_{tC\sigma(t,C)}-a_{tC})m_{tC}\\
 &\le M_te_t+\omega_t.
\end{align*}
The lower objective is treated with the reverse inequality. Summing with $B_t$ bounds any normalized pattern-flow objective within $\mathcal E$ of that pattern's true cost difference. Weighting by $\lambda_\sigma$ and summing proves that every feasible upper objective is at most $j_++\mathcal E$ and every feasible lower objective at least $j_- -\mathcal E$. The actual-policy embeddings give the other inequalities in \eqref{eq:pattern-error89}.

When the transitions and advantages are exact on a finite state space, $\delta_t=\omega_t=0$. Each fixed pattern has a unique occupation sequence; the convex mixture of their costs has the same minimum and maximum as this finite set. A retained qualification witness therefore attains each extremal policy value. More generally the quantitative inequalities prove the stated convergence whenever $\mathcal E$ tends to zero uniformly. This is not a claim that arbitrary partitions, discontinuous deterministic transitions or incomplete pattern searches satisfy that condition.

\subsection{Strict improvement in the shared-sign example}
Let $\theta=1$ select the right-hand initial action and the free action at the right-hand successor. Let $\theta=-1$ select the left-hand initial action and the costly action at the right-hand successor. All other costs are zero. Both complete policies cost zero. In the conflict relaxation, set the right-hand initial assignment and occupation to $1/2$, and the right-hand costly assignment and occupation also to $1/2$. Put the remaining initial mass on the left-hand action; choose the free action in the left-hand successor. The forbidden pair's assignments sum to one, so its conflict inequality holds. Its cost is $1/2$. Conversely that cost cannot exceed either member of the forbidden pair, whose assignments sum to one, and thus is at most $1/2$. Pattern disaggregation has only the two zero-cost occupation sequences and gives upper bound zero. The numerical test computes the three bounds by independent path enumeration and linear programming.

\section{Proof of coupled accuracy allocation}\label{supp:coupled89}
Write $f_j(\eta)=c_j\prod_t\eta_t^{-a_{jt}}$. For any vector $v$ and positive $\eta$, direct differentiation yields
\[
 v^\top\nabla^2 f_j(\eta)v
 =f_j(\eta)\left\{
       \left(\sum_ta_{jt}\frac{v_t}{\eta_t}\right)^2
       +\sum_ta_{jt}\left(\frac{v_t}{\eta_t}\right)^2\right\}\ge0.
\]
Because every coordinate appears with positive exponent somewhere, summing these forms is strictly positive for every nonzero $v$. Hence $F$ is strictly convex on the positive orthant. Its feasible set is convex and nonempty. Every coordinate is bounded above by its cap and by $\varepsilon/w_t$. If a feasible sequence sends a coordinate to zero, a term containing that coordinate tends to infinity: all other coordinates in that term have fixed finite upper bounds. A finite sublevel set containing any feasible allocation is consequently bounded away from zero and compact. Continuity gives attainment; strict convexity gives uniqueness.

After $\eta_t=\exp(x_t)$, every nonconstant term is an exponential of an affine function. The budget becomes $\sum_tw_t\exp(x_t)\le\varepsilon$ and each cap becomes $x_t\le\log\bar\eta_t$. Both the objective and constraints are convex. This proves the alternative representation, without asserting that a floating-point optimization status is a proof.

At any feasible candidate, the supporting-hyperplane inequality is
\[
 F(\eta)\ge F(\eta^0)+g\cdot(\eta-\eta^0).
\]
If $r_t=g_t+\lambda w_t+\nu_t\ge0$, positivity of $\eta$ and its budget and caps imply
\begin{align*}
 g\cdot\eta
 &=\sum_tr_t\eta_t-\lambda\sum_tw_t\eta_t-\sum_t\nu_t\eta_t\\
 &\ge-\lambda\varepsilon-\sum_t\nu_t\bar\eta_t.
\end{align*}
Substitution proves \eqref{eq:coupled-dual89}. A sufficiently small positive vector satisfies all inequalities strictly, so Slater's condition holds. At the optimum the usual stationarity and complementary-slackness equations provide nonnegative $\lambda,\nu$ with $r=0$, $\lambda(\sum_tw_t\eta_t-\varepsilon)=0$ and $\nu_t(\eta_t-\bar\eta_t)=0$. They make the lower certificate equal to the feasible objective. Equivalently, these conditions themselves give a complete sufficiency proof by the supporting-hyperplane inequality. Since $g_t<0$, the stated maximum-ratio choice with $\nu=0$ is always a nonnegative feasible certificate.

For integer exponents the objective, derivatives, constraints and certificate at rational arguments use only rational operations. The proof does not license substituting unqualified decimal derivatives when real exponents are used; outward enclosures must then make the claimed inequalities valid. For example, a rational affine lower support with verified slope bounds on the feasible box can replace the exact tangent. The finite fixtures use integer exponents and independently recompute the gradients, so no such analytic-rounding issue is hidden in their reported gaps.

If $\eta_t^0/2<\eta_t^{\rm dyad}\le\eta_t^0$, then
\[
 f_j(\eta^{\rm dyad})
 \le2^{\sum_ta_{jt}}f_j(\eta^0)\le2^P f_j(\eta^0).
\]
Summation proves the rounding bound, while decreasing coordinates preserves both the error budget and all upper caps. Finally, the equal allocation $\eta_t=\varepsilon/T$ is feasible under the corollary's hypotheses. Since $T/\varepsilon\ge1$, every monomial is at most $C(T/\varepsilon)^P$. There are at most $KT$ of them. Adding $F_0$ proves Corollary~\ref{cor:coupled-horizon89}. Neither this feasible upper construction nor the dual lower bound on its resource model is an information-theoretic lower bound for all Bellman algorithms.
